\scanpage{10}
\mspage{- 4 -}
\begin{flushright}
n\textsuperscript{o}\handwritten{354}
\end{flushright}

5) Soient $A$ un anneau gradué de type $\Delta$, $M$ un $A$-module gradué de type $\Delta$, $\handwritten{\rho}$ un homomorphisme de $\Delta$ dans un monoïde commutatif $\Delta'$ ayant un élément neutre, tel que $\handwritten{\rho}(0)=0$. Alors $A$ est un anneau gradué de type $\Delta'$ et $M$ un $A$-module gradué de type $\Delta'$ pour les graduations de type $\Delta'$ déduites des graduations de type $\Delta$ de $A$ et $M$ par le procédé de l'exemple 2 ; cela résulte aussitôt de la relation $\handwritten{\rho}(\lambda+\lambda')=\handwritten{\rho}(\lambda)+\handwritten{\rho}(\lambda')$.

En particulier, si $\Delta=\Delta_0\times\Delta_0$, on dira que $A$ (resp. $M$) est un anneau (resp. un $A$-module) \uline{bigradué} de type $\Delta_0$ ; les graduations partielles et la graduation totale correspondantes, de type $\Delta_0$, sont alors compatibles avec la structure d'anneau de $A$ (resp. la structure de module de $M$).

6) Soient $A$ un anneau gradué de type $\Delta$, $M$ un $A$-module gradué de type $\Delta$, $\lambda_0$ un élément de $\Delta$ ; pour tout $\lambda\in\Delta$, posons $M'_\lambda=M_{\lambda+\lambda_0}$ ; comme $A_\lambda M'_\mu\subset M_{\lambda+\mu+\lambda_0}=M'_{\lambda+\mu}$, les $M'_\lambda$ forment sur $M$ une nouvelle graduation de type $\Delta$ ; le $A$-module gradué de type $\Delta$ ainsi obtenu est noté $M(\lambda_0)$ et on dit qu'il s'obtient par \uline{décalage} de $\lambda_0$ de la graduation de $M$.

*7) Soit $B$ un anneau commutatif. L'anneau de polynômes $B[X]$ à une indéterminée est gradué de type $\rlap{\ensuremath{\handwritten{\underline{\underline{\phantom{N}}}}}}N$ par les sous-groupes $BX^n$ $(n\geq 0)$ (cf. Chap. III et IV).

8) Soient $B$ un anneau commutatif, $E$ un $B$-module, $Q$ une forme quadratique sur $E$, $C(Q)$ l'algèbre de Clifford de $Q$ (cf. Chap.IX,\S{}9) ; les sous-$B$-modules $C^+(Q)$ et $C^-(Q)$ forment sur $C(Q)$ une graduation de type $\rlap{\ensuremath{\handwritten{\underline{\underline{\phantom{Z}}}}}}Z/2\rlap{\ensuremath{\handwritten{\underline{\underline{\phantom{Z}}}}}}Z$ compatible avec la structure d'anneau de $C(Q)$.*

\uline{Remarques.}- 1) Les graduations les plus souvent utilisées sont de type $\rlap{\ensuremath{\handwritten{\underline{\underline{\phantom{Z}}}}}}Z$ ou de type $\rlap{\ensuremath{\handwritten{\underline{\underline{\phantom{Z}}}}}}Z^n$ ; lorsqu'on parle d'anneaux ou de modules \uline{gradués} (resp. \uline{bigradués}, \uline{trigradués}, etc.) sans préciser de quel type, il est sous-entendu qu'il s'agit de graduations de type $\rlap{\ensuremath{\handwritten{\underline{\underline{\phantom{Z}}}}}}Z$ (resp. $\rlap{\ensuremath{\handwritten{\underline{\underline{\phantom{Z}}}}}}Z^2$, $\rlap{\ensuremath{\handwritten{\underline{\underline{\phantom{Z}}}}}}Z^3$, etc.) ; on dit aussi qu'un anneau (resp. module) gradué de type $\rlap{\ensuremath{\handwritten{\underline{\underline{\phantom{N}}}}}}N$ est un anneau (resp. module) \uline{gradué à degrés positifs}.
% archive sheet number: 9

